\section{The order-10 computational exclusion}

The strict local condition studied by Kobayashi and Nakamura requires every
pair of factors in a $1$-factorization of $K_{n,n}$ to form only $4$-cycles;
such a $4$-semiregular factorization exists precisely in power-of-two
orders~\cite{KobayashiNakamura1994}.  In one Latin-square view, that condition
forces every induced two-line cycle to have length $2$.  Our F condition is
weaker: induced cycles may have any even length, and FFF requires the weaker
condition simultaneously in all three coordinate views.  The nearby strict
classification nevertheless makes the following conjecture natural.

Equivalently, view a Latin square as a triangle decomposition of
$K_{n,n,n}$.  The links of the vertices in each part form a 1-factorization
between the other two parts.  A two-line permutation cycle of length $\ell$
corresponds to an alternating link-union cycle of length $2\ell$.  Thus FFF
asks that all pair-union components in all three coupled link factorizations
have length divisible by four.  Moreover, the total number of 4-cycles in
each of the three link families is the same: in every view it is exactly the
number of intercalates of the Latin square.  This identity is a baseline
three-view coupling, although it is not by itself an obstruction.

\begin{conjecture}[Power-of-two FFF conjecture]\label{conj:power-two}
For $n>1$, an FFF Latin square of order $n$ exists if and only if $n$ is a
power of two.
\end{conjecture}

The positive direction follows from Corollary~\ref{cor:group-fff}.  The
negative direction is proved here for odd $n$ by
Proposition~\ref{prop:odd-order} and computed for $n=6$ by
Computational Theorem~\ref{cthm:small-orders}.  Even non-power-of-two orders,
including $6$ and $10$, are excluded by the computational theorems below.
The first undecided even non-power-of-two order is $12$.

Corollary~\ref{cor:2p-simple} was the first structural reduction for order
$10$: every order-$10$
FFF candidate must be congruence-simple in every isotope.  Thus central,
binary-extension, and other congruence-imprimitive constructions cannot
produce a counterexample.  This reduction alone is structural rather than an
existence decision; the complete palette computation below closes the case.

Computational Theorem~\ref{cthm:primitive-view} sharpens this again at the
level of each line action: every normalized view must generate $A_{10}$ or
$S_{10}$, and Corollary~\ref{cor:at-least-one-s10} forces at least one
$S_{10}$ view.  This is a separate necessary-condition reduction, not an
assumption used to discard candidates in the two master computations below.
Those computations cover their full normalized permutation families without
a primitive-group filter.

There is also an exact cycle-type reduction.  The seven possible cycle types
of a fixed-point-free permutation of degree $10$ with no odd cycle are
\[
 \mathcal T=\{10,8+2,6+4,6+2+2,4+4+2,
                 4+2+2+2,2+2+2+2+2\}.
\]
For one view, call the set of types occurring among its $45$ relative
two-line permutations its \emph{exact palette}.

\begin{computationaltheorem}[Order-$10$ cycle-type palette exclusions]
\label{cthm:order10-palettes}
Every view of a hypothetical order-$10$ FFF Latin square has an exact palette
containing at least six members of $\mathcal T$.  None of the $21$ five-member
subsets of $\mathcal T$ can occur as the exact palette of any view:
\begingroup
\small
\setlength{\jot}{1pt}
\begin{align*}
 &\{10,2^5,4+2+2+2,4+4+2,6+2+2\},\\
 &\{10,2^5,4+2+2+2,4+4+2,6+4\},\\
 &\{10,2^5,4+2+2+2,4+4+2,8+2\},\\
 &\{10,2^5,4+2+2+2,6+2+2,6+4\},\\
 &\{10,2^5,4+2+2+2,6+2+2,8+2\},\\
 &\{10,2^5,4+2+2+2,6+4,8+2\},\\
 &\{10,2^5,4+4+2,6+2+2,8+2\},\\
 &\{10,2^5,4+4+2,6+4,8+2\},\\
 &\{10,2^5,6+2+2,6+4,8+2\},\\
 &\{10,4+2+2+2,4+4+2,6+2+2,6+4\},\\
 &\{10,4+2+2+2,4+4+2,6+2+2,8+2\},\\
 &\{10,4+2+2+2,4+4+2,6+4,8+2\},\\
 &\{10,4+2+2+2,6+2+2,6+4,8+2\},\\
 &\{10,4+4+2,6+2+2,6+4,8+2\},\\
 &\{10,2^5,4+4+2,6+2+2,6+4\},\\
 &\{2^5,4+2+2+2,4+4+2,6+2+2,6+4\},\\
 &\{2^5,4+2+2+2,4+4+2,6+2+2,8+2\},\\
 &\{2^5,4+2+2+2,4+4+2,6+4,8+2\},\\
 &\{2^5,4+2+2+2,6+2+2,6+4,8+2\},\\
 &\{2^5,4+4+2,6+2+2,6+4,8+2\},\\
 &\{4+2+2+2,4+4+2,6+2+2,6+4,8+2\}.
\end{align*}
\endgroup
Here $2^5$ abbreviates $2+2+2+2+2$.
\end{computationaltheorem}

The lower bound combines complete exclusions of every exact palette of size
at most four.  Each of the twenty-one additional exclusions uses a complete
root-orbit cover, a fully audited compatibility graph, reduced and
unreduced searches, and an independent clique check.  Several compatibility
graphs contain Latin completions: f18 and f19 each have $1920$, f20 has
$23040$, including $19200$ with the exact f20 row palette, f21 has $1920$
with no exact f21 row palette, and f02 has $101760$, including $99840$ with
the exact f02 row palette, while f07 has $1920$ with no exact f07 row palette.
The f03 graph has $768$ cliques, split as $384$ FFT and $384$ FTF, with no
exact f03 row palette; the f04 graph has no bucket-transversal $8$-clique.
The f05 graph has $8832$ exact-f05 tables, split as $3840$ FTT, $2496$ FTF,
and $2496$ FFT.  The f11 graph has $384$ exact-f11 tables, all with pattern
FTT.  The f06 graph has $2880$ candidate tables, none with the exact f06 row
palette, split as $1152$ FTT, $864$ FTF, and $864$ FFT.  The f12 graph has
$1920$ candidate tables, of which $1152$ have the exact f12 row palette; in
aggregate they split as $1152$ FTT, $384$ FTF, and $384$ FFT.
The f08 graph has $21120$ candidate tables, none with the exact f08 row
palette; they split as $7680$ FTT, $6720$ FTF, and $6720$ FFT.  The f13
graph has $17280$ candidate tables, of which $13632$ have the exact f13 row
palette; they split as $9600$ FTT, $3840$ FTF, and $3840$ FFT.  Direct
three-view validation, not pairwise compatibility alone, excludes FFF.  The
f14 graph has $21312$ candidate tables, of which $18432$ have the exact f14
row palette; they split as $13056$ FTT, $4128$ FTF, and $4128$ FFT, with zero
FFF.  The f09 graph has $23040$ candidate tables, of which $21120$ have the
exact f09 row palette; they split equally as $7680$ FTT, $7680$ FTF, and
$7680$ FFT, with zero FFF.  The f10 graph has $55680$ candidate tables, of
which $34560$ have the exact f10 row palette; they split as $26880$ FTT,
$14400$ FTF, and $14400$ FFT, with zero FFF.  Only f15 among the exact
five-type palettes remained at that stage. The f15 graph has $32896$
candidate tables, including $13504$ with the exact f15 row palette; they
split as $6208$ FTT, $13344$ FTF, and $13344$ FFT, with zero FFF. Thus every
exact five-type palette is excluded.

The remaining palettes reduce to two disjoint normalized master families.
If the involution type $2^5$ occurs, one such line pair can be normalized and
the master permits all seven types.  This covers the six exact six-type
palettes containing $2^5$ and the exact seven-type palette.  If $2^5$ does not
occur, the lower bound in Computational
Theorem~\ref{cthm:order10-palettes} forces the unique six-type palette
\[
 \{10,4+2+2+2,4+4+2,6+2+2,6+4,8+2\},
\]
and an occurring $4+2+2+2$ pair is normalized instead.  Simultaneous
conjugation preserves all three F conditions, and exact residual-group orbit
covers reduce the two searches from $18300$ and $23760$ root tasks to $386$
and $506$ representatives.

\begin{computationaltheorem}[Nonexistence at order $10$]
\label{cthm:order10-nonexistence}
There is no FFF Latin square of order $10$.
\end{computationaltheorem}

For the noninvolution master, a complete semantic audit checks all
$14855277528$ unordered pairs among $172368$ candidates and reconstructs
$1650487056$ graph edges with zero errors.  Two complete decompositions of
the $386$ orbit representatives agree exactly at $2514104197$ search nodes
and $1769153203$ cross-view prunes, with no FFF completion.  For the
involution/all-seven master, the corresponding figures are $190080$
candidates, $18065108160$ audited pairs, $2050329600$ edges, $506$ orbit
representatives, $5001728282$ nodes, and $3586770865$ prunes, again with no
FFF completion.  The orbit identities cover every root task, and both package
validators and fail-closed mutation audits pass.  The two decompositions in
each package use the same frozen search binary and control slicing and merge
determinism; independent model checks are supplied by separate inventory
implementations, complete graph auditors, orbit reconstruction, and positive
order-$8$ controls.

By the lower bound, every hypothetical FFF view lies in exactly one of these
two masters.  Their complete exclusion proves the theorem.  This confirms
Conjecture~\ref{conj:power-two} at order $10$, not in general.

\paragraph{Completeness bridge.}
The two-master reduction fixes the identity row and a canonical second row
$\rho$: an involution of type $2^5$, or a permutation of type $4+2+2+2$.
Every further row is a permutation of ten symbols.  The candidate list
enumerates all such permutations, retaining exactly those whose relative
row types lie in the master palette and whose three-row partial table has
no already closed odd column or symbol cycle.  Candidate vertices are
adjacent exactly when their relative row type is allowed and their joint
four-row partial table passes the same cross-view check.  The eight parts
are the possible first-column images $2,\ldots,9$.  A completion therefore
selects one candidate from each part, with pairwise graph compatibility.

Write $q_r(c)=L(r,c)$ for the permutation of row $r$.  For columns $c\ne d$,
each chosen row contributes the partial map on symbols
$L(r,c)\mapsto L(r,d)$.  For symbols $u\ne v$, it contributes the
partial map on columns $q_r^{-1}(u)\mapsto q_r^{-1}(v)$.  Latin compatibility
makes every such partial map injective.  A closed odd cycle already present
in one map cannot be broken by adding rows: all its incoming and outgoing
edges are fixed.  Rejecting it is consequently sound.  At completion, the
column maps are conjugate to the usual two-column permutations and the
symbol maps are the usual two-symbol permutations, so these tests are exact.

The search branches over every remaining candidate in a selected part and
intersects the remaining parts with its graph neighbours.  An empty part
or a closed odd cross-view cycle cannot contain a completion.  Induction on
the number of unfilled parts shows that all other transversal choices are
visited.  At depth eight, the two fixed rows and eight selected rows form
a Latin square, with every row, column and symbol test enforced.  Finally,
the recorded residual conjugation actions preserve these tests and the
chosen root part; their complete orbit covers account for every initial
task.  This establishes the mathematical reduction from full FFF to the
recorded exhaustive searches.  Executable correctness and completed coverage
remain computational evidence obligations, not consequences of the
candidate counts alone.

\section{Discussion and open problems}

Order $10$ is computationally excluded by
Computational Theorem~\ref{cthm:order10-nonexistence}; a conceptual,
order-independent explanation is still sought.  The next undecided even
non-power-of-two order is $12$.  The order-$10$ encodings, proposed next
steps and diagnostic failures discussed below are historical routes and
possible sources of general principles, not a reopened existence decision.

Computational Theorem~\ref{cthm:affine-profile-saturation} also rules out a
simple continuation of the aggregate-count approach: at order $8$, over four
small prime fields, the complete affine relation space contains nothing
beyond fixed-point-free partition identities, intercalate equality, and the
global sign relation.  The same baseline admits the formal order-$10$
profile in which every line pair is a $10$-cycle.
Proposition~\ref{prop:line-sign-cut} now shows that this profile is not
realizable even in one view: all $45$ relative permutations would be odd,
whereas the odd pairs must form a cut.  The formal profile remains useful
only as evidence that aggregate affine congruences do not see this
obstruction.  A successful general obstruction must therefore retain pair
incidence or use nonlinear compatibility among the three link
factorizations.

There is a canonical cell-local refinement.  For a cell $t$, let
$a^V_\ell(t)$ count the other lines in view $V$ whose induced cycle through
$t$ has length $\ell$.  Then $\sum_\ell a^V_\ell(t)=n-1$, and, pointwise,
\[
 a^R_2(t)=a^C_2(t)=a^S_2(t),
\]
because all three quantities count the intercalates containing $t$.  Hence
the cross tensors
$M^{VW}_{\ell m}=\sum_t a^V_\ell(t)a^W_m(t)$ have exact cycle-count
marginals and retain the alignment of the three views at the same cells.
This is strictly finer than the aggregate data above.  A complete order-$8$
census gives $226$ such signatures among the $230$ FFF main classes.  Adding
weighted overlaps between individual cycle-component supports gives $229$
deterministic $1$-WL fingerprints.  The remaining pair contains one
group-isotopic and one non-group-isotopic square, and even their audited
triple-overlap and full cell/cycle-atom $1$-WL fingerprints agree.  Thus this
  finite invariant is strong but does not determine algebraic closure.  Before
the master exclusion, this pointed to sign-cut-compatible mixed profiles
rather than the already impossible all-$10$ profile; it remains useful for
possible order-independent arguments.

The flag surface of Section~\ref{sec:flag-surface} supplies such a genuinely
three-view refinement.  In the complete order-$8$ FFF census, adding only
$\operatorname{Fix}(\beta^2)$ to the cell-local signature separates all 230
main classes.  This does not classify another order: the unresolved task is
to convert the exact link, Gram, or mixed-word constraints into an
order-independent obstruction for sign-cut-compatible degree-$10$ profiles.
The exact atom projection of Theorem~\ref{thm:flag-atom-projection} removes
the large cell Gram inverse, but a verbatim order-$10$ support encoding would
already need at least $19500$ binary support entries, versus $2350$ variables
in the direct colouring model.  The next viable step must therefore eliminate
the support matrix through a genuinely compressed rank or minor consequence;
renaming all support bits is not a smaller model.
Moreover, an exact abstract construction at the order-$10$ totals satisfies
the flag count, all-even atom lengths, intercalate equality, three valid sign
cuts, Euler parity, and the componentwise atom-Gram rank simultaneously.
It deliberately decouples the formal line-pair labels from the flag
incidence and lacks the common $100$-cell support and Schur realization.
Thus no contradiction can use only those aggregate layers; a stronger
condition may retain pair-label/flag-incidence compatibility or genuine
cell-coordinate compatibility.

Theorem~\ref{thm:labelled-pair-gram} now makes the first of these intermediate
layers exact.  Its three ternary pair matrices retain every labelled
two-factor and all common flag-label codegrees.  The nontrivial Gram data
live only on three copies of the cycle space of $K_n$.  However, projecting
these labels through the coarse flag-by-cell incidence adds no condition:
the resulting Schur matrix is positive semidefinite exactly when the common
pair Gram is.  The pair-label triangle trace is highly discriminating at
order $8$---together with the cell-local signature it separates all 230 FFF
classes---but this finite classification is not an order-$10$ obstruction.
The next compression must therefore retain finer common cell support than
$B^{\mathsf T}Z$, or expose another nonlinear compatibility among the three
edge-lift sums.  Repeating the coarse Schur projection cannot succeed.

Proposition~\ref{prop:pair-label-tensor} retains the full sparse common
triple data, and Theorem~\ref{thm:endpoint-characterization} now closes its
logical realizability problem: the fibre laws and the three point-induced
endpoint edge-lift families already force sharp transitivity, one common
Latin table, and the original tensor.  Together with link divisibility this
is an exact abstract characterization of FFF Latin tensors.
Proposition~\ref{prop:intrinsic-edge-lift} makes the remaining lift condition
intrinsic: except for the explicit order-$4$ complement coset, bijectivity
and preservation of pair-label intersection recognize exactly the point
lifts.  The order-$4$ boundary also shows that local fibre and completion
geometry without global star information is insufficient.

Theorem~\ref{thm:pair-tensor-smith} gives a matrix-valued local
characterization inside this common tensor.  Odd cycles are exactly the
$\mathbb Z/2$ torsion factors of its integral slices, and FFF is equivalent
to total unimodularity of every normalized slice.  This identifies the exact
minor obstruction rather than another scalar aggregate.  It still acts one
slice at a time: a successful order-$10$ argument must couple these balanced
slices through their shared tensor entries or endpoint contractions.  The
observed global flattening ranks in the tracked partial corpus are therefore
reported only as exploratory data.

Theorem~\ref{thm:triple-endpoint-parity} supplies the first explicit
contraction that couples all slices on actual coordinate triples.  Its exact
Latin-cell values recover intercalate degrees, while all one-dimensional
binary fibres have even weight.  This closes the simplest common-cell
parity calculation, but not the order-$10$ problem: the endpoint map is onto
the even-sum point space, and the tracked partial corpus already attains mode
rank $n-1$.  A useful obstruction must therefore combine these cell values
with the balanced-slice minors of
Theorem~\ref{thm:pair-tensor-smith} or with nonlinear endpoint-lift
compatibility; line parity or flattening rank alone cannot provide the
required contradiction.

Theorem~\ref{thm:colour-parity-holonomy} completes the proposed gluing step
at the topological level.  It turns every local alternating slice
bipartition from Theorem~\ref{thm:pair-tensor-smith} into one global
$\mathbb F_2^3$ cocycle, identifies the sum of the
three colour classes with $w_1$, and forces equal quotient-colour fibres.
For order $n\equiv2\pmod4$, at least one component must have nontrivial
holonomy.  This is a genuine global restriction, but it is still equivalent
to the FFF condition rather than a contradiction.  The order-$8$ holonomy
profile is also not complete: it has $152$ values on $230$ FFF main classes.
Theorem~\ref{thm:holonomy-cell-lift} now performs the first such coupling:
on every intercalate component the three exact holonomy potentials push
forward to oriented Latin-cell vectors whose common absolute support is
exactly the multiplicity-four layer of
Theorem~\ref{thm:triple-endpoint-parity}.  This completely accounts for the
$Y$-attachment, but leaves the multiplicity-one layer $X$ untouched.  The
next obstruction must therefore couple those oriented $Y$ columns to $X$, or
transport the potentials through the endpoint lifts of
Theorem~\ref{thm:endpoint-characterization}; component topology
and the intercalate layer alone do not decide order $10$.

Theorem~\ref{thm:exact-mask-cell-pushforward} extends the cell construction
to every component, including the multiplicity-one surface: each set bit of
an exact colour mask forces the corresponding complete family of cell
marginals to vanish.  This closes the simplest linear extension of
Theorem~\ref{thm:holonomy-cell-lift}.
It still acts one component and one exact mask at a time.  A decisive
obstruction must relate different components, force additional masks to be
exact, or impose a genuinely joint equation on the $X$ and $Y$ layers.

Theorem~\ref{thm:cellwise-holonomy-fourier} now gives the complete coupling
among all exact masks at one component-cell incidence: their pushforwards are
the Fourier transform of nonnegative integral flag-label multiplicities.
Thus the separate exact-mask vectors obey exact inverse-cone and
quadratic-energy
constraints.  Gauge choices and disconnectedness still leave different
components independent.  The next missing relation must glue their label
systems through shared Latin cells or through
Theorem~\ref{thm:endpoint-characterization}.

Theorem~\ref{thm:cell-role-affine-transport} also retains all three cell
roles.  It forces equal role totals and local six-divisibility whenever the
orientation displacement is nonzero, in particular on orientable components.
Theorem~\ref{thm:cross-component-cell-parity} now glues the resulting parity
supports across all components sharing a cell and identifies their
nonintercalate sum with the multiplicity-one endpoint parity from
Theorem~\ref{thm:triple-endpoint-parity}.  This closes
the binary support layer.  It still forgets the nonzero Fourier labels,
integer masses, character signs, and relative gauges.  At this stage the
missing step was an integral or signed lift of this cover, or a canonical
comparison through the endpoint maps of
Theorem~\ref{thm:endpoint-characterization}.

Theorem~\ref{thm:gauge-invariant-cell-lift} closes the first of these routes
at the unanchored linear level: independent component gauges forbid genuine
linear signed cancellation, and odd-weight exact masks vanish pointwise.  Its
canonical cubic invariant is integral, but the explicit order-$4$
counterexample shows that the naive global cubic sum is not determined by
the intercalate degree.  Repeating unanchored sign sums cannot succeed.

Theorem~\ref{thm:quadratic-gram-gauge-obstruction} closes the proposed
endpoint-origin alternative as a universal invariant: a Klein-four
autotopism stabilizes one component while complementing two of its exact-mask
potentials, even though the full endpoint data of
Theorem~\ref{thm:endpoint-characterization} are present.  The
off-diagonal Fourier/Gram kernel is the canonical quadratic replacement and
lifts the pairwise support overlaps of
Theorem~\ref{thm:cross-component-cell-parity} integrally.  It misses
components whose
exact-mask space has dimension zero; the frozen controls contain $179048$
odd component-cell incidences of precisely this kind.
Theorem~\ref{thm:component-cell-edge-partition} now reaches that blind layer
without gauges: every component supplies three actual simple cell graphs
with a common degree vector, full line-clique edge partitions, and integral
Gram/Laplacian identities.  This also exposes the next boundary.  Separate
line graphicality does not force the three graphs to admit compatible
rainbow flag triangles or a common Latin-square realization.  A further
obstruction must therefore constrain this simultaneous cross-view
realization, not select gauges or merely repeat degree-sequence tests.
Theorem~\ref{thm:cell-companion-surface} supplies the first such coupling:
the rainbow triangles force three-periodic local links and an orientable
normalized surface with an exact genus formula.  What remains is to connect
this cell-companion topology back to the atom links of
Theorem~\ref{thm:flag-surface} or the pair-label tensor and endpoint lifts of
Proposition~\ref{prop:pair-label-tensor} and
Theorem~\ref{thm:endpoint-characterization} strongly enough to characterize,
or obstruct, a common Latin
realization.
Theorem~\ref{thm:normalized-cell-rank} now determines the complete linear
kernel of its normalized cell incidence and identifies the component
cell-edge Gram from Theorem~\ref{thm:component-cell-edge-partition} as its
coarse positive-semidefinite aggregation.  Nontrivial values $q_F=2,3$ show
that this refinement is not vacuous, but rank and positivity alone still do
not couple the normalized sheets to the atom or endpoint labels.  The next
useful target is therefore a mixed atom/cell rank, Schur, or endpoint
compatibility relation, not another isolated incidence moment.
Theorem~\ref{thm:joint-atom-cell-rank} carries out that mixed-rank step and
shows that the overlap can be highly nonconstant: one frozen order-$8$
component has twelve independent excess common face modes.  The new gap is
therefore structural rather than numerical.
Theorem~\ref{thm:common-mode-gradient} now gives the missing local
description: the common modes are exactly potentials whose three
opposite-value cochains are gradients on connected atom-transition graphs,
and it separates coarse cell modes from a normalization-only quotient.
Theorem~\ref{thm:global-common-mode} also resolves the global linear coupling:
all line sums agree, the centered modes lie in the rational Latin-trade
space, and the atom-only operator $\mathcal D_L$ gives the exact remaining
Schur nullity.  Computational
Theorem~\ref{cthm:order6-common-mode-lattice} resolves the first integral
order-$6$ layer after distinguishing occupied-cell dependencies from
classical Latin bitrades: $180$ complete-census exceptions carry a genuine
index-two common-mode quotient beyond all balanced $4+4$ exchanges.  The
intrinsic parity description is supplied by Computational
Theorem~\ref{cthm:order6-intrinsic-exchange-parity}, which also shows that the
$180$ exceptions form one precise TTT, non-group-isotopic main class.  The
complete order-$8$ FFF controls in Computational
Theorem~\ref{cthm:order8-exchange-parity-controls} close that proposed route:
FFF realizes $q_4=0,1,2,3$, including $39$ representatives with no binary
exchange quotient.  Any further obstruction must therefore retain endpoint,
cycle, or flag compatibility discarded by the bare exchange lattice, rather
than merely refine its quotient dimension.  The observed maximum $s_F=13$, the
corpus identity $p_F=q_F$, the unique observed sheet-only component, and the
baseline value on the partial tracked order-$10$ corpus are not universal
bounds or an order-$10$ obstruction.

The remaining issue is computational rather than logical.  For fixed pair
marginals only $mn^2$ tensor positions can be nonzero, at most $4500$ when
$n=10$, while a realized tensor has at most $900$ nonzero entries.  A dense
$m^3$ translation nevertheless loses this sparsity, and recognizing all edge
lifts still carries the information of the Latin table.  A useful
proof-producing implementation must store fibre types and completions
sparsely and demonstrate a genuine propagation or size advantage over the
direct colouring encoding, rather than merely rename its variables.

We tested the transparent version of that proposal.  In the three conjugate
Latin tables, two diagonal-equality bits determine each fibre type and all
completion labels, so the resulting record layer is a definitional extension
of the direct colouring model.  With $m=\binom n2$ it adds exactly $15m^2$
variables and $(18n+30)m^2$ clauses.  At order $10$ this changes the direct
balanced/gauged model from $2350$ variables and $340655$ clauses to $32725$
variables and $765905$ clauses.  Certified order-$6$ and fixed order-$8$
gates confirm semantic correctness, but the size comparison stops this
particular architecture before an order-$10$ run.  This is not a lower bound
for non-definitional eliminations; it narrows the next target to a genuinely
new compatibility relation among the three primitive line families.

The transparent intrinsic edge-map translation is not that relation.  At
order $10$, beyond the common edge-map variables and input one-hot clauses,
the existing explicit point-lift layer adds $3000$ variables and $149100$
clauses.  A direct one-hot encoding of edge-map bijectivity and all
intersection mismatches adds no point variables but $28553850$ clauses.
This is not a lower bound for optimized extended formulations, but it rules
out the naive intrinsic CNF as the next computation.

Corollary~\ref{cor:prime-congruence-series} proves the conjecture for every
quasigroup admitting a prime-factor congruence series.  Thus any counterexample
in a non-power-of-two order must lie outside this congruence-decomposable
class.  Proposition~\ref{prop:q-obstruction} gives a second reduction: at
order $10$, a QQQ candidate can fail FFF only through cycle type $5+3+2$.

A useful exact encoding follows from a basic cycle-coloring fact.  A
fixed-point-free permutation has only even cycles if and only if its points
admit a $2$-coloring that toggles along every permutation edge.  Applying one
such coloring to every line pair in every view yields a finite SAT/CP model
for FFF existence.  The repository records complete proofs of this encoding,
symmetry-breaking lemmas, and independently checked order-$6$ DRAT traces.
For order $10$, however, the complete direct SAT/CP models and all rigorously
complete symmetry splits attempted in this encoding time out.  Partial models
satisfy at most $131$ of the $135$ line-pair conditions in the best
independently validated incumbent, but that number is heuristic evidence, not
an impossibility proof.  These timeout runs are not used in Computational
Theorem~\ref{cthm:order10-nonexistence}; that theorem is proved instead by the
two complete cycle-palette master searches above.
An exact local follow-up excludes every reduced FFF table within cell-Hamming
distance $20$ of that incumbent, with an independently checked DRAT trace.
Thus no reduced FFF table lies in that radius-$20$ neighbourhood of the
incumbent.  This local diagnostic is compatible with, but logically unnecessary
for, the global nonexistence theorem and supplies no lower bound between
isotopy classes.
Residual column symmetry can be compressed to twelve fixed-row formulas by
column lex leaders.  A further exact tie-breaker under the stabilizer of the
chosen first column also passes the order-$6$ proof gates and an order-$8$
positive gate, but enlarges the tested formulas and gives no measured
order-$10$ speedup.  We therefore do not treat longer repetitions of that
encoding as evidence.

Theorem~\ref{thm:odd-subset-collision} gives an exact alternative that removes
the colouring variables.  At order $10$ it forbids collisions in the induced
actions on $3$-subsets and complementary-paired $5$-subsets.  Its complete
order-$6$ negative and fixed order-$8$ positive/negative gates are
proof-checked, but a bounded fixed order-$10$ split times out in both the
plain and redundant-colour formulations.  The theorem is therefore a new
structural characterization, not an order-$10$ decision or a demonstrated
solver speedup.

Scalar representation theory also reaches a precise boundary.  In the normal
Cayley graph on $S_{10}$ whose connection classes have all-even cycle type,
the standard character gives $|S|\leq10$, exactly the derangement bound
already attained by every Latin line family.  The full scalar-character audit
finds an exact F-view profile that remains feasible after imposing the known
three-view sign-parity and equal-intercalate aggregates.  Hence repeating
scalar Delsarte bounds or traces could not settle order $10$ at this level of
coupling.  A general representation-theoretic route must retain
matrix-valued or common-cell information across the three views.

Four concrete problems remain:
\begin{enumerate}
\item Decide Conjecture~\ref{conj:power-two} at order $12$, or prove a uniform
obstruction for even non-powers of two.
\item Find a structural explanation of the order-$10$ master exclusions that
couples the three views without exhaustive graph search.
\item Determine whether the palette lower bound and two-root reduction have
useful analogues in order $2p$ for odd primes $p>5$.
\item Determine whether the direct-product construction yields distinct
isotopy or main classes under additional hypotheses.  No such injectivity is
asserted in this paper.
\end{enumerate}

The first two questions concern existence of local odd-cycle obstructions.
They should not be conflated with the global Alon--Tarsi conjecture, whose
resolution would require substantially more than the existence or absence of
this particular local move~\cite{AlonTarsi1992}.
